% =====================================================================
%  Seksjon: Implementering av tårnmodene (side-to-side + fore-aft) i den
%           forenklede vindturbinmodellen (WindTurbineTower).
%  Kilde:   src/tops_openfast/dyn_models/windturbine_tower.py
%  Preamble: amsmath, amssymb, graphicx, booktabs, siunitx, hyperref
%           + \DeclareSIUnit{\pu}{pu}   (per-unit, used by \SI{}{\pu})
% =====================================================================

\section{Adding the Tower Structural Modes to the Reduced Wind-Turbine Model}
\label{sec:simplified-tower-model}

The OpenFAST co-simulation of \S\ref{sec:openfast-excitation} and
\S\ref{sec:tower-ss-resonance} established, in a high-fidelity aeroelastic model,
that the tower side-to-side (SS) mode can be driven by a generator-torque
oscillation. That model is, however, expensive and its electrical$\to$structure
path is partly obstructed by ROSCO (\S\ref{subsec:ss-torque-path}). To study the
interaction \emph{systematically} --- swept over frequency, contrasted against
the fore-aft (FA) mode, and combined with realistic grid events on the full
LEOGO network --- a fast, fully electrical-to-mechanical model is needed. This
section describes such a model: \texttt{WindTurbineTower}, a superset of the
reduced two-mass turbine of \S\ref{sec:frequency-analysis} that adds the first
tower side-to-side and first tower fore-aft bending modes as modal degrees of
freedom.

\subsection{Modal (assumed-mode) representation}
\label{subsec:stm-modal}

Each tower mode is represented as a single-degree-of-freedom modal oscillator,
exactly as ElastoDyn/FAST represents tower bending through a Rayleigh--Ritz mode
shape. The side-to-side displacement \(q_{\mathrm{ss}}\) obeys
\begin{equation}
  \ddot{q}_{\mathrm{ss}}
    + 2\zeta_{\mathrm{ss}}\omega_{\mathrm{ss}}\,\dot{q}_{\mathrm{ss}}
    + \omega_{\mathrm{ss}}^2\,q_{\mathrm{ss}}
  = g_{\mathrm{ss}}\,T_e,
  \qquad \omega_{\mathrm{ss}} = 2\pi f_{\mathrm{ss}},
  \label{eq:stm-ss}
\end{equation}
and the tower-top side-to-side \emph{acceleration} (the analogue of the OpenFAST
output \texttt{YawBrTAyp}) is its second derivative
\begin{equation}
  a_{\mathrm{ss}} = \ddot{q}_{\mathrm{ss}}
    = g_{\mathrm{ss}}\,T_e
      - 2\zeta_{\mathrm{ss}}\omega_{\mathrm{ss}}\,\dot{q}_{\mathrm{ss}}
      - \omega_{\mathrm{ss}}^2\,q_{\mathrm{ss}}.
  \label{eq:stm-ss-accel}
\end{equation}
Here \(T_e = P_e/(\omega_{e,\mathrm{filt}}\,\eta)\) is the same electromagnetic
torque that drives the drivetrain equation \eqref{eq:Te}, so the driving term of
\eqref{eq:stm-ss} carries the grid disturbance directly. The physical mechanism
is the nacelle reaction torque: the air-gap torque \(T_e\) acts about the
low-speed-shaft axis, and by Newton's third law the nacelle feels the equal and
opposite reaction, which rolls it laterally and bends the tower top side to side.
This is the same mechanism exploited by active side-to-side tower-damping
controllers and by ROSCO's Tower Resonance Avoidance logic, both of which
modulate generator torque to influence tower motion.

The fore-aft displacement \(q_{\mathrm{fa}}\) mirrors \eqref{eq:stm-ss} but is
driven by the rotor \emph{thrust} rather than the torque,
\begin{equation}
  \ddot{q}_{\mathrm{fa}}
    + 2\zeta_{\mathrm{fa}}\omega_{\mathrm{fa}}\,\dot{q}_{\mathrm{fa}}
    + \omega_{\mathrm{fa}}^2\,q_{\mathrm{fa}}
  = g_{\mathrm{fa}}\,\hat{F}_T,
  \qquad
  \hat{F}_T = \frac{C_T(\lambda,\beta)\,v^2}
                   {\big(C_T(\lambda,\beta)\,v^2\big)_0},
  \label{eq:stm-fa}
\end{equation}
where \(\hat{F}_T\) is the thrust normalised by its equilibrium value (so a
steady thrust produces no ring), \(\lambda = \omega_m R/v\) is the tip-speed
ratio and \(\beta\) the blade pitch. The two oscillators add four states
\(\{q_{\mathrm{ss}},\dot{q}_{\mathrm{ss}},q_{\mathrm{fa}},\dot{q}_{\mathrm{fa}}\}\)
to the model.

\subsection{Forward coupling: the tower does not react back}
\label{subsec:stm-forward}

The two modal oscillators are driven by the drivetrain/aerodynamic state
(\(T_e\) for SS, \(\hat{F}_T\) for FA) but do \emph{not} feed back onto the
drivetrain or the electrical network: they are pure output DOFs. This
one-directional (forward) coupling is deliberate and has two consequences. First,
it is physically reasonable here --- the tower-top lateral acceleration is orders
of magnitude too small to modify the air-gap torque or the delivered power.
Second, and more usefully for the study, it means that switching a mode on or off
(\(\texttt{ss\_enable}\), \(\texttt{fa\_enable}\)) leaves the electrical solution
\emph{bit-for-bit identical}. Any two runs that differ only in the tower DOFs
therefore share exactly the same grid trajectory, so a comparison isolates the
mechanical response without confounding electrical differences. The same holds
for on/off comparisons of the frequency-support droop (\S\ref{sec:outer-droop}),
which is why the droop trade-off of \S\ref{subsec:nt-tradeoff} is clean.

\subsection{Calibration against the OpenFAST-FMU}
\label{subsec:stm-calibration}

The modal frequencies and damping ratios are taken from, and the coupling gains
calibrated against, the high-fidelity OpenFAST model, so that the reduced model
reproduces the FMU tower response quantitatively rather than only qualitatively.

\paragraph{Frequencies.}
The IEA 15\,MW monopile first tower modes sit at
\(f_{\mathrm{ss}} = \SI{0.234}{\hertz}\) (side-to-side) and
\(f_{\mathrm{fa}} = \SI{0.235}{\hertz}\) (fore-aft); the two are nearly
degenerate in frequency, which makes the \emph{drive path}, not the frequency,
the discriminator between them (\S\ref{sec:shaft-excitation}).

\paragraph{Side-to-side gain \(g_{\mathrm{ss}}\).}
For the oscillator \eqref{eq:stm-ss} driven at its natural frequency, the
steady-state acceleration amplitude follows from the transfer function of
\eqref{eq:stm-ss-accel} at \(\omega=\omega_{\mathrm{ss}}\),
\begin{equation}
  \big|a_{\mathrm{ss}}(\omega_{\mathrm{ss}})\big|
  = \frac{g_{\mathrm{ss}}}{2\zeta_{\mathrm{ss}}}\,\big|\Delta T_e\big|.
  \label{eq:stm-gss}
\end{equation}
The OpenFAST-FMU torque sweep (\S\ref{sec:tower-ss-resonance}) measured a peak
side-to-side acceleration of \(\approx\SI{0.259}{\metre\per\second\squared}\)
(\texttt{YawBrTAyp} at \SI{0.234}{\hertz}) for a \(\pm\SI{10}{\percent}\)
generator-torque modulation (\(|\Delta T_e|\approx\SI{0.10}{\pu}\)) at
\(\zeta_{\mathrm{ss}}\approx\num{0.05}\). Solving \eqref{eq:stm-gss} gives a
first estimate \(g_{\mathrm{ss}}\approx\SI{0.26}{\metre\per\second\squared\per\pu}\),
which is then refined by an equivalent reduced-model torque sweep
(\texttt{sweep\_ss\_reduced.py}) that injects the same \(\pm\SI{10}{\percent}\)
modulation on \(T_e\) and matches the FMU peak, yielding the calibrated default
\begin{equation}
  g_{\mathrm{ss}} = \SI{0.2684}{\metre\per\second\squared\per\pu},
  \qquad
  \zeta_{\mathrm{ss}} = \num{0.05}.
  \label{eq:stm-gss-value}
\end{equation}

\paragraph{Fore-aft gain and effective damping.}
The fore-aft gain \(g_{\mathrm{fa}} = \SI{0.102}{\metre\per\second\squared\per\pu}\)
and effective damping \(\zeta_{\mathrm{fa}} = \num{0.0125}\) are calibrated the
same way, from the FMU fore-aft ring-down and forced-response at
\(v = \SI{11}{\metre\per\second}\). The damping is termed \emph{effective}
because the strong aerodynamic damping of the fore-aft mode --- when the tower
top moves downwind the rotor sees less wind, the thrust drops and opposes the
motion --- is folded into a single modal \(\zeta_{\mathrm{fa}}\) rather than
resolved explicitly. Notably \(\zeta_{\mathrm{fa}}\) here is \emph{small}: the
reason the fore-aft mode is not the resonance-prone one is not heavy damping but
the weak thrust drive path (\S\ref{sec:shaft-excitation}).

Table~\ref{tab:stm-params} collects the calibrated modal parameters. With them
the reduced model reproduces the FMU side-to-side resonance peak to within a few
percent (\S\ref{sec:new-tests}) while running in real time against the full
LEOGO network.

\begin{table}[t]
  \centering
  \caption{Calibrated tower modal parameters of \texttt{WindTurbineTower},
  taken from / matched to the OpenFAST-FMU IEA 15\,MW monopile model.}
  \label{tab:stm-params}
  \begin{tabular}{lccc}
    \toprule
    Mode & \(f\) [\si{\hertz}] & \(\zeta\) [--] &
      \(g\) [\si{\metre\per\second\squared\per\pu}] \\
    \midrule
    Side-to-side (torque-driven) & 0.234 & 0.050 & 0.2684 \\
    Fore-aft (thrust-driven)     & 0.235 & 0.0125 & 0.102 \\
    \bottomrule
  \end{tabular}
\end{table}

\subsection{Ring-down time constant}
\label{subsec:stm-tau}

A consequence of the light side-to-side damping worth stating explicitly, since
it recurs in the scenarios of \S\ref{sec:new-tests}, is the modal decay time
constant
\begin{equation}
  \tau_{\mathrm{ss}} = \frac{1}{\zeta_{\mathrm{ss}}\,\omega_{\mathrm{ss}}}
    = \frac{1}{0.05\cdot 2\pi\cdot 0.234}
    \approx \SI{13.6}{\second}.
  \label{eq:stm-tau}
\end{equation}
A broadband ``kick'' (a load step) therefore excites a free side-to-side
ring-down that decays with \(\tau_{\mathrm{ss}}\approx\SI{13.6}{\second}\),
whereas a sustained sinusoidal drive at \(f_{\mathrm{ss}}\) builds the mode up
over \(\sim 3\text{--}4\,\tau_{\mathrm{ss}}\approx\SI{45}{\second}\) to a large
steady amplitude (\S\ref{subsec:nt-resonance}). The distinction between the two
--- ring-\emph{down} versus build-\emph{up} --- is the clearest experimental
signature of resonant versus broadband excitation.

\subsection*{Sources}
\label{subsec:stm-sources}

\begin{itemize}
  \item Modal (assumed-mode / Rayleigh--Ritz) tower-bending DOFs, as used by
        ElastoDyn/OpenFAST: Jonkman (2007), \emph{Dynamics Modeling and Loads
        Analysis of an Offshore Floating Wind Turbine}, NREL/TP-500-41958,
        \href{https://www.nrel.gov/docs/fy08osti/41958.pdf}{nrel.gov/docs/fy08osti/41958.pdf};
        Jonkman \& Buhl (2005), \emph{FAST User's Guide}, NREL/EL-500-38230,
        \href{https://www.nrel.gov/docs/fy06osti/38230.pdf}{nrel.gov/docs/fy06osti/38230.pdf};
        OpenFAST documentation, \href{https://openfast.readthedocs.io/}{openfast.readthedocs.io}.
  \item IEA \SI{15}{\mega\watt} reference turbine (modal frequencies, rotor and
        tower data): Gaertner et al.\ (2020), \emph{Definition of the IEA
        15-Megawatt Offshore Reference Wind Turbine}, NREL/TP-5000-75698,
        \href{https://www.nrel.gov/docs/fy20osti/75698.pdf}{nrel.gov/docs/fy20osti/75698.pdf};
        model repository
        \href{https://github.com/IEAWindTask37/IEA-15-240-RWT}{github.com/IEAWindTask37/IEA-15-240-RWT}.
  \item Fore-aft aerodynamic tower damping (the ``aerodynamic damper''):
        K\"uhn (2001), \emph{Dynamics and Design Optimisation of Offshore Wind
        Energy Conversion Systems}, PhD thesis, TU Delft; Salzmann \& van der
        Tempel (2005), \emph{Aerodynamic damping in the design of support
        structures for offshore wind turbines}, Copenhagen Offshore Wind
        Conference.
\end{itemize}
